\section{Deeper Physical Applications}
\label{sec:ch3-physical-applications}

The same four ideas recur throughout theoretical physics: a density, a current or flux, a local differential law, and a global integral statement. This section develops that architecture before later chapters attach the full physical dynamics.

\subsection{Fluid dynamics: local expansion and rotation}
For a velocity field $\mathbf v$, the divergence measures local volume expansion and the vorticity $\boldsymbol\omega=\nabla\times\mathbf v$ measures local rotational tendency. Conservation of mass is
\[
\frac{\partial\rho}{\partial t}+\nabla\cdot(\rho\mathbf v)=0.
\]
For constant density this reduces to $\nabla\cdot\mathbf v=0$, the incompressibility condition. It does not imply zero curl: shear and vortex flows can be incompressible.

\begin{physical}
The differential equation describes what happens inside an infinitesimal control volume. Gauss' theorem converts it into a statement that the rate of mass loss inside a finite region equals the outward mass flux across its boundary.
\end{physical}

\subsection{Electrostatics and gravitation: sources and flux}
Electrostatic and Newtonian gravitational fields share inverse-square geometry. In differential form,
\[
\nabla\cdot\mathbf E=\frac{\rho_q}{\varepsilon_0},
\qquad
\nabla\cdot\mathbf g=-4\pi G\rho_m.
\]
The signs encode the physical distinction between positive electric sources and universally attractive gravitational mass. In source-free regions, potentials satisfy Laplace's equation. Boundary values then control the interior field through uniqueness principles.

\subsection{Magnetism: circulation without monopole sources}
Writing $\mathbf B=\nabla\times\mathbf A$ automatically gives
\[
\nabla\cdot\mathbf B=0.
\]
Thus magnetic flux lines do not begin or end at ordinary isolated sources. Stokes' theorem relates the circulation of the vector potential around a boundary to magnetic flux through a spanning surface:
\[
\oint_{\partial S}\mathbf A\cdot d\mathbf l=\iint_S\mathbf B\cdot d\mathbf S.
\]
This relation later becomes central in electromagnetic induction, gauge theory, Berry phase, and the quantum Hall effect.

\subsection{Diffusion, waves, and spectral structure}
The Laplacian distinguishes three important dynamical patterns:
\[
\partial_tu=D\nabla^2u,
\qquad
\partial_t^2u=c^2\nabla^2u,
\qquad
-\nabla^2\psi=\lambda\psi.
\]
They represent diffusion, waves, and spatial eigenmodes. The operator is the same, but the time law and boundary conditions produce fundamentally different behavior: irreversible smoothing, oscillatory propagation, or discrete mode spectra.

\subsection{Quantum probability conservation}
A wavefunction carries probability density $\rho=|\psi|^2$. Schrödinger evolution produces a probability current $\mathbf j$ satisfying
\[
\frac{\partial |\psi|^2}{\partial t}+\nabla\cdot\mathbf j=0.
\]
The equation has exactly the same local-conservation structure as fluid mass and electric charge. Quantum mechanics changes what the density and current mean, not the underlying geometry of conservation.

\begin{roadmapbox}[One mathematical architecture]
\[
\text{local source}\;\longleftrightarrow\;\text{divergence},\qquad
\text{local rotation}\;\longleftrightarrow\;\text{curl},
\]
\[
\text{boundary flux}\;\longleftrightarrow\;\text{Gauss},\qquad
\text{boundary circulation}\;\longleftrightarrow\;\text{Stokes}.
\]
This architecture is reused in continuum mechanics, Maxwell theory, gauge theory, and quantum transport.
\end{roadmapbox}

\subsection{From Chapter 3 to the rest of the book}
Chapter~4 uses vector-valued curves to describe trajectories. Newtonian mechanics then interprets gradients as forces derived from potentials. Electromagnetism combines divergence and curl into a closed field theory. Quantum mechanics promotes the Laplacian to a kinetic-energy operator and conservation laws to probability flow. The reader should therefore regard vector calculus not as a preliminary language, but as part of the conceptual structure of physics itself.
